\chapter{Combinatorics, Probability and Games}\label{ch:combinatorics}

Counting problems ask ``in how many ways?'', usually modulo a prime, and their answers are formulas assembled from a few building
blocks: binomial coefficients, inclusion--exclusion, symmetry, and a handful of famous sequences. This chapter builds that toolbox,
adds the algebra that turns hard counts into easy ones (Stirling numbers, polynomial interpolation, determinants), and ends with two
neighbouring subjects: expected values, and games in which two players alternate moves.

\begin{chapterintro}{\silver\ \gold\ \plat\ \ioi}{Chapter~\ref{ch:numbertheory} (modular inverses) and
Section~\ref{sec:probdp} and Section~\ref{sec:gamedp} (probability and game DPs).}
  \item count with the product, sum and complement rules, bijections and double counting;
  \item compute binomial coefficients modulo a prime in $\Oh{1}$, and modulo small primes with Lucas's theorem;
  \item count lattice paths, distributions and bracket sequences with stars and bars and the reflection principle;
  \item subtract overcounts with inclusion--exclusion, and count up to symmetry with Burnside's lemma;
  \item use Stirling numbers, Lagrange interpolation, and determinants (matrix-tree theorem, LGV lemma);
  \item compute expected values by linearity, and win impartial games with Nim and Sprague--Grundy theory.
\end{chapterintro}

% ======================================================================================================================
\section[Counting principles]{Counting principles \hfill \silver}\label{sec:counting}

Almost every count is built from five principles.
\begin{itemize}
  \item \textbf{Product rule}: a sequence of $k$ independent choices with $c_1, \dots, c_k$ options gives $c_1 c_2 \cdots c_k$ outcomes.
        ``Independent'' means that the number of options at each step does not depend on the earlier choices, even if the options
        themselves do.
  \item \textbf{Sum rule}: if the outcomes split into disjoint cases, add the counts of the cases.
  \item \textbf{Complement}: count everything and subtract the outcomes you do not want; ``at least one'' is usually ``all minus none''.
  \item \textbf{Bijection}: count a different set that is in one-to-one correspondence with the one you want, such as subsets and
        bit strings, or lattice paths and bracket sequences.
  \item \textbf{Double counting}: count the same quantity in two ways, typically as a sum over objects of their parts, and as a sum
        over parts of the objects that contain them. The ``contribution technique'' of Chapter~\ref{ch:patterns} is this principle.
\end{itemize}

Choosing $k$ elements out of $n$ comes in four flavours (Table~\ref{tab:selections}), depending on whether order matters and whether
an element may be chosen more than once.

\begin{table}[htbp]
\centering\small
\begin{tabular}{@{}lll@{}}
\toprule
choose $k$ of $n$ & order matters & order does not matter \\
\midrule
without repetition & $n^{\underline{k}} = n (n-1) \cdots (n - k + 1) = \frac{n!}{(n-k)!}$ & $\binom{n}{k} = \frac{n!}{k!\,(n-k)!}$ \\[2pt]
with repetition & $n^k$ & $\binom{n + k - 1}{k}$ (stars and bars) \\
\bottomrule
\end{tabular}
\caption{The four ways to choose $k$ elements from $n$.}
\label{tab:selections}
\end{table}

When the objects are not all different, divide out the orderings that cannot be told apart: a word with $n$ letters in which letter $i$
occurs $c_i$ times has $\frac{n!}{c_1!\, c_2! \cdots}$ distinct arrangements, the \defn{multinomial coefficient}.

\begin{problem}{Creating Strings II (CSES)}
How many distinct strings can be formed from the letters of a given string of length $n \le 10^6$? Answer modulo $10^9 + 7$.
\solution The multinomial coefficient $n! / \prod_c (\text{count of } c)!$, with the divisions done by multiplying with precomputed
inverse factorials (next section). The program takes 0.09\,s for $n = 10^6$.
\end{problem}

\begin{code}
void solve() {  // nCr (cpncr) is pasted above solve()
    string s;
    cin >> s;
    init_comb(sz(s));
    vi cnt(26);
    for (char c : s) cnt[c - 'a']++;
    ll ans = fact[sz(s)];  // all orders, then forget the order of equal letters
    for (ll c : cnt) ans = ans * inv_fact[c] % MOD;
    print(ans);
}
\end{code}

% ======================================================================================================================
\section[Binomial coefficients]{Binomial coefficients \hfill \gold}\label{sec:binomial}

$\binom{n}{k}$ counts the $k$-element subsets of an $n$-element set. Modulo a prime $p > n$, precompute $i!$ for all $i \le N$ and then
the inverse factorials: $(N!)^{-1}$ with one fast exponentiation, and $(i - 1)!^{-1} = i!^{-1} \cdot i$ going down. After this $\Oh{N}$
preparation, every $\binom{n}{k} = n! \cdot (k!)^{-1} \cdot ((n-k)!)^{-1}$ costs two multiplications.

\usetemplate{nCr / factorials, Pascal}{Ctrl+Alt+T, N / Ctrl+Alt+Y, Shift+N}{cpncr, cppascal}

\begin{problem}{Binomial Coefficients (CSES)}
Answer $10^5$ queries $\binom{a}{b} \bmod (10^9 + 7)$ with $0 \le b \le a \le 10^6$.
\solution Factorials and inverse factorials up to $10^6$, once; every query is then $\Oh{1}$. 0.12\,s in total.
\end{problem}

\begin{code}
void solve() {  // nCr (cpncr) is pasted above solve()
    init_comb(1000000);  // factorials and inverse factorials, once
    int q;
    cin >> q;
    while (q--) {
        int a, b;
        cin >> a >> b;
        print(nCr(a, b));
    }
}
\end{code}

\paragraph{Pascal's rule.} $\binom{n}{k} = \binom{n-1}{k-1} + \binom{n-1}{k}$: a $k$-subset of $\{1, \dots, n\}$ either contains $n$ or
not. The rule needs no division, so Pascal's triangle computes binomial coefficients modulo \emph{any} number, including composite
moduli where inverses do not exist, in $\Oh{n^2}$ (Figure~\ref{fig:pascal}). Exact values fit in 64 bits up to $\binom{66}{33}
\approx 7.2 \cdot 10^{18}$.

\begin{figure}[htbp]
\centering
\begin{tikzpicture}
  \begin{scope}
    \tikzset{pn/.style={vtx, minimum size=6.6mm, font=\sffamily\scriptsize}}
    \foreach \rowv [count=\rn from 0] in {{1}, {1,1}, {1,2,1}, {1,3,3,1}, {1,4,6,4,1}, {1,5,10,10,5,1}, {1,6,15,20,15,6,1}} {
      \foreach \vv [count=\kk from 0] in \rowv \node[pn] (C\rn-\kk) at (0.78*\kk - 0.39*\rn, -0.68*\rn) {\vv};
    }
    \node[pn, cool] at (C3-1) {3};
    \node[pn, cool] at (C3-2) {3};
    \node[pn, hot] at (C4-2) {6};
    \draw[->, line width=0.8pt, draw=fgblue, shorten >=1pt, shorten <=1pt] (C3-1) -- (C4-2);
    \draw[->, line width=0.8pt, draw=fgblue, shorten >=1pt, shorten <=1pt] (C3-2) -- (C4-2);
    \node[title] at (0, 0.75) {Pascal's triangle};
  \end{scope}
  \begin{scope}[shift={(7.4, 0)}, x=0.27cm, y=0.27cm]
    \foreach \rn in {0, ..., 15} \foreach \kk in {0, ..., \rn}
      \node[draw=fgline!80, line width=0.3pt, minimum size=2.5mm, inner sep=0pt] at (\kk - \rn / 2, -\rn) {};
    \foreach \rn/\kk in {0/0, 1/0, 1/1, 2/0, 2/2, 3/0, 3/1, 3/2, 3/3, 4/0, 4/4, 5/0, 5/1, 5/4, 5/5, 6/0, 6/2, 6/4, 6/6, 7/0,
        7/1, 7/2, 7/3, 7/4, 7/5, 7/6, 7/7, 8/0, 8/8, 9/0, 9/1, 9/8, 9/9, 10/0, 10/2, 10/8, 10/10, 11/0, 11/1, 11/2, 11/3, 11/8,
        11/9, 11/10, 11/11, 12/0, 12/4, 12/8, 12/12, 13/0, 13/1, 13/4, 13/5, 13/8, 13/9, 13/12, 13/13, 14/0, 14/2, 14/4, 14/6,
        14/8, 14/10, 14/12, 14/14, 15/0, 15/1, 15/2, 15/3, 15/4, 15/5, 15/6, 15/7, 15/8, 15/9, 15/10, 15/11, 15/12, 15/13, 15/14,
        15/15}
      \node[draw=fgblue, fill=fgblue!65, line width=0.3pt, minimum size=2.5mm, inner sep=0pt] at (\kk - \rn / 2, -\rn) {};
    \node[title] at (0, 2.8) {odd entries, rows 0 to 15};
  \end{scope}
\end{tikzpicture}
\caption{Left: each entry is the sum of the two above it, $\binom42 = \binom31 + \binom32$. Right: the odd entries of the first 16 rows
form a Sierpi\'nski triangle. By Lucas's theorem, $\binom{n}{k}$ is odd exactly when every 1-bit of $k$ is also a 1-bit of $n$: 81 of
the 136 entries.}
\label{fig:pascal}
\end{figure}

\begin{fact}[Identities worth knowing]
Each has a one-line counting proof.
\begin{itemize}
  \item $\binom{n}{k} = \binom{n}{n-k}$ (choose the elements to leave out) and $\sum_k \binom{n}{k} = 2^n$ (all subsets).
  \item $k \binom{n}{k} = n \binom{n-1}{k-1}$ (choose a $k$-subset and a leader in it, or the leader first).
  \item $\sum_{i=k}^{n} \binom{i}{k} = \binom{n+1}{k+1}$, the hockey stick (classify the $(k+1)$-subsets of $\{0, \dots, n\}$ by their
        largest element $i$).
  \item $\sum_k \binom{m}{k}\binom{n}{r-k} = \binom{m+n}{r}$, Vandermonde (choose $r$ from two groups).
  \item Monotone lattice paths from $(0, 0)$ to $(a, b)$: $\binom{a+b}{a}$ (choose which steps go right).
\end{itemize}
\end{fact}

\begin{trick}[Lucas's and Kummer's theorems]
For a prime $p$ smaller than $n$, factorials vanish modulo $p$ and the inverse-factorial method breaks. Write $n$ and $k$ in base $p$;
then $\binom{n}{k} \equiv \prod_i \binom{n_i}{k_i} \pmod p$ (\textbf{Lucas}), where each small binomial is taken from a table of size
$p^2$. Kummer's theorem says more: the power of $p$ that divides $\binom{n}{k}$ is the number of carries when $k$ and $n - k$ are
added in base $p$. For $p = 2$, $\binom{n}{k}$ is odd exactly when \cc{(k & n) == k} (Figure~\ref{fig:pascal}). For composite moduli,
compute modulo each prime power and combine with the Chinese remainder theorem (Section~\ref{sec:crt}).
\end{trick}

% ======================================================================================================================
\section[Stars and bars and lattice paths]{Stars and bars and lattice paths \hfill \gold}\label{sec:starsbars}

The number of ways to put $n$ identical items into $k$ distinguishable boxes, empty boxes allowed, is $\binom{n + k - 1}{k - 1}$: write
the items as $n$ stars and separate the boxes by $k - 1$ bars; every arrangement of the $n + k - 1$ symbols is one distribution
(Figure~\ref{fig:starsbars}). If every box must get at least one item, first put one item in each box: $\binom{n - 1}{k - 1}$. The same
count answers ``how many solutions in non-negative integers does $x_1 + \dots + x_k = n$ have?''.

\begin{figure}[htbp]
\centering
\begin{tikzpicture}
  \foreach \xx/\kind in {0/s, 1/s, 2/b, 3/s, 4/b, 5/b, 6/s, 7/s, 8/s} {
    \if\kind s
      \node[font=\Large, text=fgblue] at (0.62*\xx, 0) {$\star$};
    \else
      \draw[line width=1.6pt, draw=fgverm] (0.62*\xx, -0.32) -- (0.62*\xx, 0.32);
    \fi
  }
  \foreach \xa/\xb/\cnt in {0/1/2, 3/3/1, 4.5/4.5/0, 6/8/3}
    \draw[mbrace] (0.62*\xa - 0.25, -0.42) -- node[below=6pt, lbl] {\cnt} (0.62*\xb + 0.25, -0.42);
  \node[notel, anchor=west, text width=7.2cm] at (6.1, -0.1) {6 stars and 3 bars: the boxes receive 2, 1, 0 and 3 items. Choosing
    which 3 of the 9 positions hold bars gives $\binom{9}{3} = 84$ distributions of 6 items into 4 boxes.};
\end{tikzpicture}
\caption{Stars and bars: a distribution of identical items is an arrangement of stars and bars.}
\label{fig:starsbars}
\end{figure}

\begin{problem}{Distributing Apples (CSES)}
In how many ways can $m$ identical apples be given to $n$ children ($n, m \le 10^6$)?
\solution $\binom{n + m - 1}{n - 1}$, in 0.09\,s including the factorial table.
\end{problem}

\begin{code}
void solve() {  // nCr (cpncr) is pasted above solve()
    int n, m;
    cin >> n >> m;
    init_comb(n + m);
    print(nCr(n + m - 1, n - 1));  // m stars and n - 1 bars
}
\end{code}

\paragraph{Paths that avoid forbidden cells.} On a grid, the monotone paths from $(y_1, x_1)$ to $(y_2, x_2)$ number
$\binom{(y_2 - y_1) + (x_2 - x_1)}{y_2 - y_1}$. With a few forbidden cells, inclusion--exclusion becomes manageable once every bad path is
charged to the \emph{first} forbidden cell it visits (Figure~\ref{fig:gridtraps}).

\begin{figure}[htbp]
\centering
\begin{tikzpicture}[scale=0.62]
  \foreach \yy in {1, ..., 6} \foreach \xx in {1, ..., 6}
    \node[draw=fgline, line width=0.5pt, minimum size=6.2mm, inner sep=0pt, fill=white] at (\xx, -\yy) {};
  \foreach \yy/\xx in {2/4, 4/2, 4/5}
    \node[draw=fgverm, line width=1pt, minimum size=6.2mm, inner sep=0pt, fill=fgverm!18] at (\xx, -\yy) {};
  \draw[line width=2pt, draw=fgverm, line cap=round] (1, -1) -- (1, -3) -- (2, -3) -- (2, -4);
  \draw[line width=2pt, draw=fgblue, line cap=round] (2, -4) -- (3, -4) -- (3, -5) -- (4, -5) -- (4, -6) -- (6, -6);
  \fill[fgink] (1, -1) circle (3pt);
  \fill[fgink] (6, -6) circle (3pt);
  \node[lbl, anchor=east] at (0.55, -1) {start};
  \node[lbl, anchor=west] at (6.45, -6) {end};
  \node[notel, anchor=west, text width=7.4cm] at (7.6, -3.5) {Traps (red cells) at $(2, 4)$, $(4, 2)$ and $(4, 5)$. The bad path
    shown is counted once, at its first trap $(4, 2)$: its red part is one of the $f(4, 2)$ paths from the start that reach $(4, 2)$
    without touching an earlier trap, and its blue part is any of the $\binom{2 + 4}{2}$ paths from $(4, 2)$ to the end.};
\end{tikzpicture}
\caption{Grid Paths II: every path through a trap is charged to the first trap it visits.}
\label{fig:gridtraps}
\end{figure}

\begin{problem}{Grid Paths II (CSES)}
Count the monotone paths from the top-left to the bottom-right corner of an $n \times n$ grid ($n \le 10^6$) that avoid $m \le 1000$
traps, modulo $10^9 + 7$.
\solution Sort the traps; a trap can be reached only through traps that come before it. Let $f_i$ be the number of paths from the
start to trap $i$ that avoid all earlier traps: all paths to it, minus, for every earlier trap $j$, the $f_j$ paths that reach $j$ first
times the paths from $j$ to $i$. Treating the target as one more trap, its $f$ is the answer. $\Oh{m^2}$ binomials after an $\Oh{n}$
table: 0.10\,s.
\end{problem}

\begin{code}
void solve() {  // nCr (cpncr) is pasted above solve()
    int n, m;
    cin >> n >> m;
    vector<pii> t(m);
    for (auto& [y, x] : t) cin >> y >> x;
    sort(all(t));  // a trap can only be reached through traps before it
    t.push_back({n, n});  // the target, treated as the last trap
    init_comb(2 * n);
    auto paths = [&](pii a, pii b) {  // monotone paths from a to b
        ll dy = b.first - a.first, dx = b.second - a.second;
        return dy < 0 || dx < 0 ? 0LL : nCr(dy + dx, dy);
    };
    vi f(m + 1);  // f[i]: paths from (1, 1) to t[i] avoiding earlier traps
    rep(i, m + 1) {
        f[i] = paths({1, 1}, t[i]);
        rep(j, i) f[i] = (f[i] - f[j] * paths(t[j], t[i]) % MOD + MOD) % MOD;
    }
    print(f[m]);
}
\end{code}

% ======================================================================================================================
\section[Catalan numbers and the reflection principle]{Catalan numbers and the reflection principle \hfill \plat}\label{sec:catalan}

Read a bracket sequence as a walk: ``('' is a step up, ``)'' a step down. The sequence is balanced exactly when the walk ends at height 0
and never goes below 0. How many walks of $2n$ steps from 0 to 0 never go below 0? All walks from 0 to 0 number $\binom{2n}{n}$. Take a
bad one, find the first time it touches $-1$, and reflect the part before that moment in the line $-1$: the result is a walk from $-2$
to 0. The reflection is reversible, so bad walks correspond one to one to walks from $-2$ to 0, which number $\binom{2n}{n+1}$
(Figure~\ref{fig:ballot}). The good walks number
\[
  C_n = \binom{2n}{n} - \binom{2n}{n+1} = \frac{1}{n+1}\binom{2n}{n}, \qquad 1, 1, 2, 5, 14, 42, 132, 429, 1430, 4862, \dots
\]
the \defn{Catalan numbers}.

\begin{figure}[htbp]
\centering
\begin{tikzpicture}[x=0.62cm, y=0.5cm]
  \fill[fgverm!8] (0, -1) rectangle (10, -3.4);
  \draw[line width=0.4pt, draw=fgline] (0, -3) grid (10, 2);
  \draw[line width=0.8pt, draw=fgink!70] (0, 0) -- (10, 0);
  \draw[line width=0.8pt, draw=fgverm, dashed] (0, -1) -- (10, -1);
  \node[lbl, anchor=east] at (-0.15, 0) {0};
  \node[lbl, anchor=east] at (-0.15, -1) {$-1$};
  \node[lbl, anchor=east] at (-0.15, -2) {$-2$};
  \draw[line width=1.8pt, draw=fgblue, line cap=round, line join=round] (0, 0) -- (1, 1) -- (2, 0) -- (3, -1) -- (4, 0) -- (5, 1)
    -- (6, 0) -- (7, 1) -- (8, 2) -- (9, 1) -- (10, 0);
  \draw[line width=1.8pt, draw=fgverm, line cap=round, line join=round, dash pattern=on 5pt off 3pt] (0, -2) -- (1, -3) -- (2, -2)
    -- (3, -1);
  \fill[fgverm] (3, -1) circle (3.5pt);
  \node[font=\sffamily\scriptsize, text=fgverm, anchor=north west] at (3.1, -1.1) {first touch of $-1$};
  \node[notel, anchor=west, text width=5.6cm] at (10.6, -0.5) {The walk \texttt{())(()(())} (blue) goes below 0. Reflecting its
    first three steps in the line $-1$ gives a walk from $-2$ to 0 (red, then blue). Of the 252 walks from 0 to 0 in 10 steps, 210
    are bad, leaving $C_5 = 42$ balanced sequences.};
\end{tikzpicture}
\caption{The reflection principle: bad walks from 0 correspond to walks from $-2$.}
\label{fig:ballot}
\end{figure}

The same walks count binary trees with $n$ nodes, triangulations of a convex $(n+2)$-gon, ways to multiply $n + 1$ matrices, monotone
lattice paths below the diagonal of an $n \times n$ grid, and non-crossing matchings of $2n$ points on a circle (Figure~\ref{fig:catalan}).
Each satisfies the recurrence $C_{n+1} = \sum_{i=0}^{n} C_i C_{n-i}$: split a balanced sequence at the bracket that closes its first
``('', a binary tree at its root, a triangulation at the triangle on one fixed side.

\begin{figure}[htbp]
\centering
\begin{tikzpicture}
  \foreach \brk [count=\ii from 0] in {((())), (()()), (())(), ()(()), ()()()}
    \node[font=\ttfamily\small] at (2.3*\ii, 1.25) {\brk};
  \foreach \ii in {0, ..., 4} {
    \begin{scope}[shift={(2.3*\ii, 0)}]
      \foreach \vv in {0, ..., 4} \coordinate (P\ii-\vv) at (90 + 72*\vv:0.62);
      \draw[line width=0.8pt, draw=fgink!80, fill=fgfill] (P\ii-0) -- (P\ii-1) -- (P\ii-2) -- (P\ii-3) -- (P\ii-4) -- cycle;
      \pgfmathtruncatemacro{\da}{mod(\ii + 2, 5)}
      \pgfmathtruncatemacro{\db}{mod(\ii + 3, 5)}
      \draw[line width=1.2pt, draw=fgverm] (P\ii-\ii) -- (P\ii-\da) (P\ii-\ii) -- (P\ii-\db);
      \fill[fgverm] (P\ii-\ii) circle (1.8pt);
    \end{scope}
  }
\end{tikzpicture}
\caption{$C_3 = 5$: the five balanced sequences of three bracket pairs, and the five triangulations of a pentagon.}
\label{fig:catalan}
\end{figure}

\begin{problem}{Bracket Sequences II (CSES)}
How many balanced bracket sequences of length $n \le 10^6$ begin with a given prefix?
\solution Scan the prefix; if its height ever drops below 0, the answer is 0. Otherwise it ends at height $b$, and the remaining $r$
brackets form a walk from $b$ to 0 that never goes below 0. Among the $r$ steps, $d = (r + b)/2$ go down, so there are $\binom{r}{d}$
walks in all; reflecting the part before the first touch of $-1$ maps the bad ones to walks from $-b - 2$ to 0, which have $d + 1$ down
steps. The answer is $\binom{r}{d} - \binom{r}{d+1}$. 0.07\,s.
\end{problem}

\begin{code}
void solve() {  // nCr (cpncr) is pasted above solve()
    int n;
    string s;
    cin >> n >> s;
    int b = 0;  // the balance: open minus closed brackets so far
    for (char c : s) {
        b += c == '(' ? 1 : -1;
        if (b < 0) {  // a prefix with more ')' than '(' never recovers
            print(0);
            return;
        }
    }
    int r = n - sz(s);  // the rest is a path of r steps from height b to 0
    if (r < b || (r - b) % 2) {
        print(0);
        return;
    }
    init_comb(r);
    int down = (r + b) / 2;  // the number of ')' among the r brackets
    print((nCr(r, down) - nCr(r, down + 1) + MOD) % MOD);  // minus bad paths
}
\end{code}

% ======================================================================================================================
\section[Inclusion--exclusion]{Inclusion--exclusion \hfill \plat}\label{sec:incexc}

To count the elements of a union, add the sets, subtract the pairwise intersections, add the triple intersections, and so on:
\[
  \Bigl| \bigcup_{i=1}^{k} A_i \Bigr| = \sum_{\emptyset \ne S \subseteq \{1, \dots, k\}} (-1)^{|S| + 1} \Bigl| \bigcap_{i \in S} A_i
  \Bigr| .
\]
An element that lies in exactly $t \ge 1$ of the sets is counted $\binom t1 - \binom t2 + \binom t3 - \dots = 1 - (1 - 1)^t = 1$ time
(Figure~\ref{fig:venn}). The method fits whenever the intersections are easy to count but the union is not: ``at least one condition
holds'', ``no forbidden event happens'' (the complement), ``every one of $k$ things is used''.

\begin{figure}[htbp]
\centering
\begin{tikzpicture}
  \begin{scope}[blend group=multiply]
    \fill[fgsky!45] (0, 0) circle (1.35);
    \fill[fgorange!40] (1.6, 0) circle (1.35);
    \fill[fggreen!35] (0.8, -1.3) circle (1.35);
  \end{scope}
  \draw[line width=0.9pt, draw=fgblue] (0, 0) circle (1.35);
  \draw[line width=0.9pt, draw=fgorange!85!black] (1.6, 0) circle (1.35);
  \draw[line width=0.9pt, draw=fggreen] (0.8, -1.3) circle (1.35);
  \node[font=\sffamily\small\bfseries, text=fgblue] at (-1.3, 1.25) {$A_2$};
  \node[font=\sffamily\small\bfseries, text=fgorange!85!black] at (2.9, 1.25) {$A_3$};
  \node[font=\sffamily\small\bfseries, text=fggreen] at (0.8, -2.95) {$A_5$};
  \foreach \px/\py/\vv in {-0.5/0.3/8, 2.1/0.3/4, 0.8/-1.95/2, 0.8/0.55/4, -0.05/-0.95/2, 1.65/-0.95/1, 0.8/-0.45/1}
    \node[font=\sffamily\small] at (\px, \py) {\vv};
  \node[notel, anchor=west, text width=7cm] at (3.6, -0.6) {The numbers up to 30 divisible by 2, 3 or 5, split by which of the
    three they are divisible by. Inclusion--exclusion: $15 + 10 + 6 - 5 - 3 - 2 + 1 = 22$; the centre number (30) is added three times,
    subtracted three times and added once. The other $30 - 22 = 8 = \varphi(30)$ numbers are coprime to 30.};
\end{tikzpicture}
\caption{Inclusion--exclusion for three sets: every region is counted exactly once.}
\label{fig:venn}
\end{figure}

\usetemplate{Inclusion--exclusion (divisible by any)}{Ctrl+Alt+Y, 3}{cpincexc}

\begin{problem}{Prime Multiples (CSES)}
Given $n \le 10^{18}$ and $k \le 20$ distinct primes, count the integers in $[1, n]$ divisible by at least one of them.
\solution For a subset of the primes, the numbers divisible by all of them are the multiples of their product, $\lfloor n /
\prod \rfloor$ of them. Sum over the $2^k$ subsets with alternating signs. Products overflow quickly, so stop multiplying as soon as the
product exceeds $n$ (the term is then 0), testing \cc{prod > n / p} before multiplying. 0.22\,s for 20 primes.
\end{problem}

\begin{code}
void solve() {  // count_divisible_by_any (cpincexc) is pasted above solve()
    ll n;
    int k;
    cin >> n >> k;
    vi a(k);
    for (ll& x : a) cin >> x;
    print(count_divisible_by_any(n, a));  // odd subsets add, even subtract
}
\end{code}

\paragraph{Derangements.} A permutation with no fixed point is a \defn{derangement}. Let $A_i$ be the permutations that fix $i$; any $j$
of these sets intersect in $(n - j)!$ permutations, so inclusion--exclusion gives $D_n = \sum_{j=0}^{n} (-1)^j \binom{n}{j} (n-j)! = n!
\sum_{j=0}^{n} \frac{(-1)^j}{j!}$, and $D_n / n! \to 1/e$. A recurrence avoids the alternating sum: in a derangement, element $n$ goes
to some position $i \ne n$ ($n - 1$ choices); either $i$ goes to $n$ (a swap, leaving $D_{n-2}$ derangements of the rest), or it does
not (then $i$ effectively takes the role of $n$, leaving $D_{n-1}$). So $D_n = (n - 1)(D_{n-1} + D_{n-2})$, with values $1, 0, 1, 2, 9,
44, 265, 1854, \dots$

\begin{problem}{Christmas Party (CSES)}
$n \le 10^6$ children each bring a gift and receive one, and no child may receive their own. Count the ways modulo $10^9 + 7$.
\solution $D_n$ by the recurrence: 0.06\,s.
\end{problem}

\begin{code}
void solve() {
    int n;
    cin >> n;
    vi d(max(n + 1, 2LL));  // d[k]: derangements of k children
    d[0] = 1, d[1] = 0;
    rep(k, 2, n + 1) d[k] = (k - 1) * (d[k - 1] + d[k - 2]) % MOD;
    print(d[n]);
}
\end{code}

\begin{fact}[Surjections]
The functions from an $n$-set onto a $k$-set (every element of the $k$-set is hit) number $\sum_{j=0}^{k} (-1)^j \binom{k}{j}
(k - j)^n$: inclusion--exclusion over the elements that are missed. This equals $k! \cdot S(n, k)$, where $S(n, k)$ is a Stirling number
of the second kind (Section~\ref{sec:stirling}).
\end{fact}

\begin{trick}[Inclusion--exclusion over subsets of a large ground set]
When the conditions are indexed by the subsets of a $k$-element set and $k \le 20$, the $2^k$ terms can be computed for all subsets at
once with the subset-sum (zeta and M\"obius) transforms of Section~\ref{sec:sosdp}: ``exactly the set $S$'' from ``at least the set $S$''
is a M\"obius transform over supersets.
\end{trick}

% ======================================================================================================================
\section[Burnside's lemma]{Burnside's lemma \hfill \ioi}\label{sec:burnside}

To count objects \emph{up to symmetry}, such as necklaces up to rotation or colourings of a grid up to turning it, list the symmetries
as a group $G$ of transformations and count the \emph{orbits}: two objects are equivalent when a symmetry maps one to the other.
\textbf{Burnside's lemma} says that the number of orbits is the average number of objects fixed by a symmetry:
\[
  \#\text{orbits} = \frac{1}{|G|} \sum_{g \in G} \mathrm{fix}(g), \qquad \mathrm{fix}(g) = \#\{x : g(x) = x\}.
\]
(Count the pairs $(g, x)$ with $g(x) = x$ in two ways. For each object $x$ there are $|G| / |\text{orbit of } x|$ such $g$, so every
orbit contributes exactly $|G|$ pairs in total.) The point is that fixed objects are easy to count: a colouring is fixed by a rotation
exactly when it is constant on every cycle of positions that the rotation moves around, so $\mathrm{fix}(g) = (\text{colours})^{\text{cycles of } g}$ (Figure~\ref{fig:necklaces}).

\begin{figure}[htbp]
\centering
\begin{tikzpicture}
  \foreach \ncol/\nrow/\ba/\bb/\bc/\bd in {0/0/0/0/0/0, 0/1/0/0/0/1, 1/1/0/0/1/0, 2/1/0/1/0/0, 3/1/1/0/0/0, 0/2/0/0/1/1,
      1/2/0/1/1/0, 2/2/1/0/0/1, 3/2/1/1/0/0, 0/3/0/1/0/1, 1/3/1/0/1/0, 0/4/0/1/1/1, 1/4/1/0/1/1, 2/4/1/1/0/1, 3/4/1/1/1/0,
      0/5/1/1/1/1} {
    \begin{scope}[shift={(1.05*\ncol, -0.82*\nrow)}]
      \draw[line width=0.5pt, draw=fgline] (0, 0) circle (0.25);
      \foreach \bit/\ang in {\ba/90, \bb/0, \bc/-90, \bd/180} {
        \ifnum\bit=1
          \fill[fgink] (\ang:0.25) circle (0.085);
        \else
          \filldraw[fill=white, draw=fgink, line width=0.5pt] (\ang:0.25) circle (0.085);
        \fi
      }
    \end{scope}
  }
  \foreach \nrow/\osz in {0/1, 1/4, 2/4, 3/2, 4/4, 5/1} \node[lbl, anchor=east] at (-0.55, -0.82*\nrow) {orbit of size \osz};
  \node[notel, anchor=west, text width=6.4cm] at (4.1, -2.05) {The $2^4 = 16$ two-coloured necklaces of 4 beads, one orbit per row.
    Rotation by 0 fixes all 16, by one bead only the 2 single-coloured necklaces, by two beads the 4 that repeat with period 2, and by
    three beads again 2. Burnside: $(16 + 2 + 4 + 2)/4 = 6$ orbits.};
\end{tikzpicture}
\caption{Necklaces of length 4 with two colours, up to rotation.}
\label{fig:necklaces}
\end{figure}

\begin{problem}{Counting Necklaces (CSES)}
Count the necklaces of $n$ beads with $m$ colours up to rotation, modulo $10^9 + 7$ ($n, m \le 10^6$).
\solution Rotation by $i$ beads splits the positions into $\gcd(i, n)$ cycles, so it fixes $m^{\gcd(i, n)}$ necklaces. The answer is
$\frac{1}{n} \sum_{i=0}^{n-1} m^{\gcd(i, n)}$; grouping the $i$ by the value of the gcd gives the classic $\frac{1}{n}\sum_{d \mid n}
\varphi(n/d)\, m^d$. 0.10\,s.
\end{problem}

\begin{code}
void solve() {  // modpow (cpmodpow) is pasted above solve()
    int n, m;
    cin >> n >> m;
    vi pw(n + 1, 1);  // pw[g] = m^g
    rep(g, 1, n + 1) pw[g] = pw[g - 1] * m % MOD;
    ll sum = 0;  // rotating by i fixes the m^gcd(i, n) colourings with period i
    rep(i, n) sum += pw[gcd(i, n)];
    print(sum % MOD * modinv(n) % MOD);  // Burnside: average over n rotations
}
\end{code}

\begin{problem}{Counting Grids (CSES)}
Count the black-and-white colourings of an $n \times n$ grid up to rotation by multiples of $90^\circ$ ($n \le 10^9$).
\solution Four rotations. The identity fixes all $2^{n^2}$ colourings. A half turn pairs up the cells, except the centre when $n$ is
odd, so it has $\lceil n^2/2 \rceil$ cycles; a quarter turn (either direction) groups them in fours, with $\lceil n^2/4 \rceil$ cycles.
The answer is $(2^{n^2} + 2 \cdot 2^{\lceil n^2/4 \rceil} + 2^{\lceil n^2/2 \rceil}) / 4$. The exponent $n^2 \le 10^{18}$ is used
directly by fast exponentiation.
\end{problem}

\begin{code}
void solve() {  // modpow (cpmodpow) is pasted above solve()
    ll n;
    cin >> n;
    ll c = n * n;  // cells; a quarter turn has (c + 3) / 4 cycles of cells
    ll fix = modpow(2, c) + 2 * modpow(2, (c + 3) / 4) + modpow(2, (c + 1) / 2);
    print(fix % MOD * modinv(4) % MOD);  // average over the four rotations
}
\end{code}

\begin{pitfall}[Rotations or also reflections?]
``Necklaces'' usually allow only rotations; ``bracelets'' may also be turned over, which doubles the group (the dihedral group of
$2n$ elements) and adds the reflections' fixed points to the sum. Read the statement carefully, and check small cases by brute force:
for $n = 4$ and two colours there are 6 necklaces and also 6 bracelets, but for $n = 6$ there are 14 necklaces and 13 bracelets.
\end{pitfall}

% ======================================================================================================================
\section[Stirling numbers]{Stirling numbers \hfill \ioi}\label{sec:stirling}

The \defn{Stirling number of the second kind} $S(n, k)$ counts the ways to partition $n$ labelled elements into $k$ non-empty unlabelled
groups. Look at element $n$: it is either alone in its group, leaving $S(n - 1, k - 1)$ partitions of the others, or it joins one of the
$k$ groups of a partition of the others, so
\[
  S(n, k) = S(n - 1, k - 1) + k \cdot S(n - 1, k), \qquad S(0, 0) = 1
\]
(Figure~\ref{fig:stirling}). The \defn{Bell number} $B_n = \sum_k S(n, k)$ counts all partitions: $1, 1, 2, 5, 15, 52, 203, 877,
\dots$ The \defn{Stirling number of the first kind} $\left[{n \atop k}\right]$ counts the permutations of $n$ elements with exactly $k$
cycles, with $\left[{n \atop k}\right] = \left[{n-1 \atop k-1}\right] + (n - 1) \left[{n-1 \atop k}\right]$: element $n$ forms its own
cycle, or is inserted after one of the $n - 1$ others in their cycles.

\begin{figure}[htbp]
\centering
\begin{tikzpicture}
  \tikzset{sc/.style={cellc, minimum width=9mm, minimum height=6.5mm, font=\sffamily\small}}
  \foreach \kk in {0, ..., 6} \node[lbl] at (0.9*\kk, 0.6) {$k = \kk$};
  \foreach \rowv [count=\rn from 0] in {{1}, {0,1}, {0,1,1}, {0,1,3,1}, {0,1,7,6,1}, {0,1,15,25,10,1}, {0,1,31,90,65,15,1}} {
    \node[lbl, anchor=east] at (-0.6, -0.68*\rn) {$n = \rn$};
    \foreach \vv [count=\kk from 0] in \rowv \node[sc] (S\rn-\kk) at (0.9*\kk, -0.68*\rn) {\vv};
  }
  \node[sc, cool] at (S4-2) {7};
  \node[sc, cool] at (S4-3) {6};
  \node[sc, hot] at (S5-3) {25};
  \node[notel, anchor=west, text width=5.8cm] at (6.5, -2.4) {$S(5, 3) = S(4, 2) + 3 \cdot S(4, 3) = 7 + 18 = 25$ (red from
    the two blue cells): element 5 is
    alone, next to one of the 7 ways to split $\{1, 2, 3, 4\}$ in two, or it joins one of the 3 groups of one of the 6 ways to
    split $\{1, 2, 3, 4\}$ in three. The row sums are the Bell numbers 1, 1, 2, 5, 15, 52, 203.};
\end{tikzpicture}
\caption{Stirling numbers of the second kind.}
\label{fig:stirling}
\end{figure}

\usetemplate{Stirling numbers (both kinds) and Bell numbers}{Ctrl+Alt+U, S}{cpstirling}

\paragraph{Powers and falling factorials.} The most useful identity converts powers into falling factorials $x^{\underline{j}} = x (x-1)
\cdots (x - j + 1)$:
\[
  x^k = \sum_{j=0}^{k} S(k, j)\, x^{\underline{j}} ,
\]
because both sides count the functions from a $k$-set to an $x$-set, the right side grouped by the partition of the $k$-set into
the preimages of the $j$ values that are hit. Falling factorials sum as easily as powers integrate: $\sum_{i=0}^{n} i^{\underline{j}} =
\frac{(n+1)^{\underline{j+1}}}{j + 1}$. Together they give the sum of $k$-th powers in $\Oh{k^2}$.

\begin{problem}{Sums of powers}
Compute $1^k + 2^k + \dots + n^k \bmod (10^9 + 7)$ for $n \le 10^{18}$ and $1 \le k \le 1000$.
\solution $\sum_{i=0}^{n} i^k = \sum_{j} S(k, j) \frac{(n+1)^{\underline{j+1}}}{j+1}$; the falling factorial is built up one factor at
a time, and $n$ enters only through these factors, so its size does not matter. 0.08\,s for $k = 1000$, most of it filling the
Stirling table.
\end{problem}

\begin{code}
void solve() {  // init_stirling (cpstirling), modinv (cpmodpow) pasted above
    ll n;
    int k;
    cin >> n >> k;
    init_stirling(k);
    ll ans = 0, fall = 1;  // fall: (n + 1) n (n - 1) ... (n + 1 - j)
    rep(j, k + 1) {
        fall = fall * ((n + 1 - j) % MOD) % MOD;
        // i^k = sum_j S2(k, j) i(i-1)...(i-j+1); summing over i = 0..n
        // turns i(i-1)...(i-j+1) into (n+1) n ... (n+1-j) / (j+1)
        ans = (ans + S2[k][j] * fall % MOD * modinv(j + 1)) % MOD;
    }
    print(ans);
}
\end{code}

% ======================================================================================================================
\section[Lagrange interpolation]{Lagrange interpolation \hfill \ioi}\label{sec:lagrange}

A polynomial of degree less than $k$ is determined by its values at $k$ points, and \textbf{Lagrange's formula} writes it down
directly:
\[
  P(x) = \sum_{i=0}^{k-1} y_i \prod_{j \ne i} \frac{x - x_j}{x_i - x_j} .
\]
Each term is $y_i$ at $x_i$ and 0 at the other points, so the sum takes the right values, and the polynomial is unique because the
difference of two solutions would be a polynomial of degree less than $k$ with $k$ roots. For general points the formula costs
$\Oh{k^2}$. For the consecutive points $0, 1, \dots, k - 1$ the denominators are $\pm\, i!\,(k - 1 - i)!$ and the numerators are a prefix
product times a suffix product, so a value at any $x$ costs $\Oh{k}$ (Figure~\ref{fig:lagrange}).

\begin{figure}[htbp]
\centering
\begin{tikzpicture}
\begin{axis}[cpplot, width=8.4cm, height=5.6cm, xmin=-0.3, xmax=4.6, ymin=-2, ymax=34, xlabel={$x$},
  ylabel={$S(x) = 1^2 + \dots + x^2$}, xtick={0, 1, 2, 3, 4}, ytick={0, 5, 14, 30}]
  \addplot[domain=0:4.4, samples=80, line width=1.2pt, draw=fgblue, no marks] {x*(x+1)*(2*x+1)/6};
  \addplot[only marks, mark=*, mark size=2.6pt, fill=fgverm, draw=fgverm] coordinates {(0,0) (1,1) (2,5) (3,14)};
  \addplot[only marks, mark=square*, mark size=2.6pt, fill=fggreen, draw=fggreen] coordinates {(4,30)};
\end{axis}
\node[notel, anchor=west, text width=5.0cm] at (8.6, 2.1) {$S(x) = \sum_{i=1}^{x} i^2$ is a polynomial of degree 3, so the four
  values $0, 1, 5, 14$ (red) determine it everywhere: interpolation predicts $S(4) = 30$ (green) and $S(10^9)$ just as easily.};
\end{tikzpicture}
\caption{Four points determine a cubic.}
\label{fig:lagrange}
\end{figure}

\usetemplate{Lagrange interpolation}{Ctrl+Alt+U, I}{cplagrange}

\begin{problem}{Sums of powers, faster}
Compute $1^k + 2^k + \dots + n^k \bmod (10^9 + 7)$ for $n \le 10^{18}$ and $k \le 10^6$.
\solution $S(x) = \sum_{i=1}^{x} i^k$ is a polynomial of degree $k + 1$ in $x$: compute its values at $x = 0, \dots, k + 1$ directly
and interpolate at $n$. $\Oh{k \log k}$ instead of the Stirling method's $\Oh{k^2}$: 0.11\,s for $k = 10^6$.
\end{problem}

\begin{code}
void solve() {  // lagrange_consecutive (cplagrange) is pasted above solve()
    ll n;
    int k;
    cin >> n >> k;
    vi y(k + 2);  // S(x) = 1^k + ... + x^k is a polynomial of degree k + 1
    rep(x, 1, k + 2) y[x] = (y[x - 1] + lag_pow(x, k)) % MOD;
    print(lagrange_consecutive(y, n));  // k + 2 values determine it everywhere
}
\end{code}

\begin{trick}[Guess that the answer is a polynomial]
Many counting answers are polynomials in a parameter, often for a structural reason (a sum of a polynomial over a range, a DP whose
transitions are polynomial in $n$), and the degree $d$ can be bounded. Then compute the first $d + 1$ answers with a slow brute force or
DP, and interpolate to the requested $n$, however large. Check the guess by predicting one extra small value.
\end{trick}

% ======================================================================================================================
\section[Determinants, spanning trees and paths]{Determinants, spanning trees and paths \hfill \ioi}\label{sec:matrixtree}

The determinant of an $n \times n$ matrix modulo a prime is computed by Gaussian elimination in $\Oh{n^3}$: bring the matrix to
triangular form with row operations, track the sign of every row swap, and multiply the diagonal. Two theorems turn determinants into
counting tools.

\paragraph{The matrix-tree theorem.} The \defn{Laplacian} of a multigraph is the matrix $L = D - A$, with the degrees on the diagonal
and minus the number of edges between $u$ and $v$ off it. \textbf{Kirchhoff's theorem}: the number of spanning trees equals the
determinant of $L$ with any one row and the same column removed (Figure~\ref{fig:matrixtree}). For the complete graph it gives
\textbf{Cayley's formula} $n^{n-2}$, which Prüfer codes (Section~\ref{sec:prufer}) prove by a bijection.

\begin{figure}[htbp]
\centering
\begin{tikzpicture}
  \tikzset{mv/.style={vtx, minimum size=6mm, font=\sffamily\scriptsize}}
  \node[mv] (G1) at (0, 1.6) {1};
  \node[mv] (G2) at (1.6, 1.6) {2};
  \node[mv] (G3) at (1.6, 0) {3};
  \node[mv] (G4) at (0, 0) {4};
  \draw[ed] (G1) -- (G2) (G2) -- (G3) (G3) -- (G4) (G4) -- (G1) (G1) -- (G3);
  \node[font=\small, anchor=west] at (2.5, 0.8) {$\det \begin{pmatrix} 3 & -1 & -1 \\ -1 & 2 & -1 \\ -1 & -1 & 3 \end{pmatrix} = 8$};
  \foreach \tr [count=\ti from 0] in {{1/2, 2/3, 3/4}, {1/2, 2/3, 4/1}, {1/2, 3/4, 4/1}, {1/2, 3/4, 1/3}, {1/2, 4/1, 1/3},
      {2/3, 3/4, 4/1}, {2/3, 3/4, 1/3}, {2/3, 4/1, 1/3}} {
    \pgfmathtruncatemacro{\tcol}{mod(\ti, 4)}
    \pgfmathtruncatemacro{\trow}{\ti / 4}
    \begin{scope}[shift={(7.1 + 1.45*\tcol, 1.55 - 1.45*\trow)}]
      \coordinate (T\ti-1) at (0, 0.85);
      \coordinate (T\ti-2) at (0.85, 0.85);
      \coordinate (T\ti-3) at (0.85, 0);
      \coordinate (T\ti-4) at (0, 0);
      \draw[line width=0.5pt, draw=fgline, dashed] (T\ti-1) -- (T\ti-2) -- (T\ti-3) -- (T\ti-4) -- cycle (T\ti-1) -- (T\ti-3);
      \foreach \ea/\eb in \tr \draw[line width=1.6pt, draw=fgblue] (T\ti-\ea) -- (T\ti-\eb);
      \foreach \vv in {1, ..., 4} \fill[fgink] (T\ti-\vv) circle (1.6pt);
    \end{scope}
  }
\end{tikzpicture}
\caption{The square 1--2--3--4 with the diagonal 1--3. Its Laplacian without the row and column of vertex 4 has determinant 8, and
indeed there are 8 spanning trees: every 3 of the 5 edges except the two triangles.}
\label{fig:matrixtree}
\end{figure}

\usetemplate{Determinant mod $p$, matrix-tree theorem}{Ctrl+Alt+U, Shift+D}{cpdet}

\begin{problem}{Counting spanning trees}
Count the spanning trees of a multigraph with $n \le 500$ vertices and $m \le 10^5$ edges, modulo $10^9 + 7$.
\solution Build the Laplacian without the last vertex and take its determinant: $\Oh{n^3}$, 0.15\,s for $n = 500$. Parallel edges
simply add up in the Laplacian, and self-loops are ignored.
\end{problem}

\begin{code}
void solve() {  // count_spanning_trees (cpdet) is pasted above solve()
    int n, m;
    cin >> n >> m;
    vector<pii> e(m);
    for (auto& [u, v] : e) cin >> u >> v, u--, v--;
    print(count_spanning_trees(n, e));  // a minor of the Laplacian (Kirchhoff)
}
\end{code}

\paragraph{The LGV lemma.} In a directed acyclic graph, let $e(a, b)$ be the number of paths from $a$ to $b$. For sources $a_1, \dots,
a_k$ and sinks $b_1, \dots, b_k$, the \textbf{Lindstr\"om--Gessel--Viennot lemma} says that $\det\bigl(e(a_i, b_j)\bigr)$ counts the
systems of $k$ vertex-disjoint paths from the sources to the sinks, each weighted by the sign of the permutation that matches sources to
sinks. In planar settings such as grids, where only the identity matching can be disjoint, the determinant simply counts the
non-intersecting path systems (Figure~\ref{fig:lgv}).

\begin{figure}[htbp]
\centering
\begin{tikzpicture}[scale=0.62]
  \draw[line width=0.4pt, draw=fgline] (0, 0) grid (4, 4);
  \draw[line width=2pt, draw=fgblue, line cap=round, line join=round] (0, 1) -- (0, 3) -- (1, 3) -- (1, 4) -- (3, 4);
  \draw[line width=2pt, draw=fgverm, line cap=round, line join=round] (1, 0) -- (3, 0) -- (3, 2) -- (4, 2) -- (4, 3);
  \foreach \px/\py/\nm/\pos in {0/1/a_1/left, 1/0/a_2/below, 3/4/b_1/above, 4/3/b_2/right}
    \node[circle, fill=fgink, inner sep=1.6pt, label={[font=\sffamily\scriptsize]\pos:$\nm$}] at (\px, \py) {};
  \node[notel, anchor=west, text width=7.6cm] at (5.2, 2) {Monotone paths from $a_1 = (0, 1)$ to $b_1 = (3, 4)$ and from $a_2 = (1,
    0)$ to $b_2 = (4, 3)$. $e(a_1, b_1) = e(a_2, b_2) = \binom63 = 20$ and $e(a_1, b_2) = e(a_2, b_1) = \binom62 = 15$, so
    $20 \cdot 20 - 15 \cdot 15 = 175$ pairs of paths share no point, as a brute force confirms.};
\end{tikzpicture}
\caption{The LGV lemma: a determinant of path counts counts non-intersecting pairs of paths.}
\label{fig:lgv}
\end{figure}

% ======================================================================================================================
\section[Probability and expected value]{Probability and expected value \hfill \gold}\label{sec:probability}

With finitely many equally likely outcomes, $P(E) = |E| / |\Omega|$, and every counting technique becomes a probability technique.
Section~\ref{sec:probdp} computed probabilities and expectations with DPs; this section collects the tools that avoid a DP altogether.

\paragraph{Linearity and indicators.} $E[X + Y] = E[X] + E[Y]$ always, even for dependent variables. Write the quantity of interest as a
sum of \emph{indicator variables}, each 1 when some event happens and 0 otherwise; the expectation is then the sum of the events'
probabilities, and each event is usually simple even when the whole is not.

\begin{fact}[Expectations by indicators]
\begin{itemize}
  \item A random permutation has on average 1 fixed point ($n$ events of probability $\frac1n$), $H_n = 1 + \frac12 + \dots + \frac1n$
        cycles, and $H_n$ left-to-right maxima (position $i$ is a maximum with probability $\frac1i$).
  \item Collecting all $n$ coupons, drawn uniformly at random, takes $n H_n \approx n \ln n$ draws on average: after $i$ distinct
        coupons, a new one appears with probability $\frac{n-i}{n}$, so it takes $\frac{n}{n-i}$ draws on average.
  \item An experiment that succeeds with probability $p$ needs $\frac1p$ tries on average.
  \item For a non-negative integer variable, $E[X] = \sum_{x \ge 1} P(X \ge x)$, the right tool for maxima: $P(\max \le x)$ is a product
        of independent probabilities (Candy Lottery, Section~\ref{sec:probdp}).
\end{itemize}
\end{fact}

\begin{problem}{Inversion Probability (CSES)}
Each $x_i$ ($n \le 100$) is chosen uniformly from $1, \dots, r_i$ ($r_i \le 100$). What is the expected number of inversions?
\solution An inversion is a pair $i < j$ with $x_i > x_j$; by linearity the answer is $\sum_{i<j} P(x_i > x_j)$. For one pair, count
the value pairs $(a, b)$ with $a > b$: for each $a$ there are $\min(a - 1, r_j)$ of them (Figure~\ref{fig:inversion}). $\Oh{n^2 r}$
operations, 0.06\,s.
\end{problem}

\begin{figure}[htbp]
\centering
\begin{tikzpicture}[scale=0.55]
  \foreach \av in {1, ..., 5} \foreach \bv in {1, ..., 3} {
    \pgfmathtruncatemacro{\inv}{\av > \bv ? 1 : 0}
    \ifnum\inv=1
      \node[draw=fgblue, fill=fgsky!40, line width=0.6pt, minimum size=5.5mm, inner sep=0pt] at (\av, \bv) {};
    \else
      \node[draw=fgline, fill=white, line width=0.6pt, minimum size=5.5mm, inner sep=0pt] at (\av, \bv) {};
    \fi
  }
  \foreach \av in {1, ..., 5} \node[lbl] at (\av, 0.2) {\av};
  \foreach \bv in {1, ..., 3} \node[lbl] at (0.2, \bv) {\bv};
  \node[lbl] at (3, -0.55) {$a = x_i$};
  \node[lbl, rotate=90] at (-0.55, 2) {$b = x_j$};
  \node[notel, anchor=west, text width=6.4cm] at (6.3, 2) {$r_i = 5$, $r_j = 3$: of the 15 equally likely value pairs, the 9
    shaded ones have $x_i > x_j$ (column $a$ contributes $\min(a - 1, 3)$ cells), so this pair is an inversion with probability
    $9/15 = 0.6$.};
\end{tikzpicture}
\caption{The probability that one pair is inverted.}
\label{fig:inversion}
\end{figure}

\begin{code}
void solve() {
    int n;
    cin >> n;
    vi r(n);
    for (ll& x : r) cin >> x;
    ld e = 0;  // linearity: add up P(x_i > x_j) over all pairs i < j
    rep(i, n) rep(j, i + 1, n) {
        ll good = 0;  // pairs (a, b) with a <= r_i, b <= r_j and a > b
        rep(a, 1, r[i] + 1) good += min(a - 1, r[j]);
        e += (ld)good / (r[i] * r[j]);
    }
    cout << fixed << setprecision(6) << e << '\n';
}
\end{code}

\begin{pitfall}[Exact answers modulo a prime]
``Print the expected value as $P \cdot Q^{-1} \bmod p$'' asks for an exact rational answer: keep every probability as a residue, divide
with modular inverses, and never pass through floating point. Conversely, problems that print decimals rarely need more than
\cc{double}; sums of many small probabilities lose less precision when added from smallest to largest.
\end{pitfall}

% ======================================================================================================================
\section[Impartial games]{Impartial games \hfill \plat}\label{sec:games}

Section~\ref{sec:gamedp} classified the positions of a single game as winning or losing with a DP. Many games are \emph{sums} of
independent games, such as several heaps of which a move changes exactly one. Their positions are too many for a DP, and the theory of
impartial games (both players have the same moves, and whoever cannot move loses) handles them.

\paragraph{Nim.} There are several heaps, and a move removes any positive number of stones from one heap. \textbf{Bouton's theorem}:
the player to move wins exactly when the XOR of the heap sizes is non-zero (Figure~\ref{fig:nim}). If the XOR is $s \ne 0$, take a heap
$h$ that has the highest bit of $s$ set; then $h \oplus s < h$, and reducing $h$ to $h \oplus s$ makes the XOR 0. If the XOR is 0, every
move changes one heap and therefore the XOR, so the opponent faces a non-zero XOR again. The final position (all heaps empty) has
XOR 0 and is lost for the player to move, so XOR 0 positions are exactly the losing ones.

\begin{figure}[htbp]
\centering
\begin{tikzpicture}
  \foreach \hx/\hv/\nm in {0/3/A, 1.4/4/B, 2.8/5/C} {
    \foreach \jj in {1, ..., \hv} \node[circle, draw=fgorange!85!black, fill=fgorange!25, line width=0.6pt, minimum size=4.6mm,
      inner sep=0pt] at (\hx, 0.5*\jj - 0.2) {};
    \node[lbl] at (\hx, -0.4) {\nm: \hv};
  }
  \node[circle, draw=fgverm, line width=1pt, minimum size=4.6mm, inner sep=0pt, dashed] at (0, 1.3) {};
  \node[circle, draw=fgverm, line width=1pt, minimum size=4.6mm, inner sep=0pt, dashed] at (0, 0.8) {};
  \begin{scope}[shift={(4.6, 2.3)}]
    \foreach \nm/\bits/\yy in {A/011/0, B/100/-0.55, C/101/-1.1} {
      \node[lbl, anchor=east] at (0, \yy) {\nm};
      \node[font=\ttfamily\small, anchor=west] at (0.1, \yy) {\bits};
    }
    \draw[line width=0.6pt, draw=fgink!60] (0.1, -1.42) -- (1.0, -1.42);
    \node[lbl, anchor=east] at (0, -1.75) {XOR};
    \node[font=\ttfamily\small\bfseries, text=fgverm, anchor=west] at (0.1, -1.75) {010};
  \end{scope}
  \node[notel, anchor=west, text width=6.2cm] at (6.2, 1.2) {Heaps 3, 4, 5 have XOR $011 \oplus 100 \oplus 101 = 010 \ne 0$: the
    player to move wins. Heap A has the bit 010 set, and $3 \oplus 2 = 1$, so take two stones from A (dashed). Now $1 \oplus 4 \oplus 5
    = 0$, and every reply breaks the zero again.};
\end{tikzpicture}
\caption{Nim: the XOR of the heap sizes decides the game.}
\label{fig:nim}
\end{figure}

\paragraph{Sprague--Grundy.} Every position $x$ of an impartial game gets a \defn{Grundy number}: $g(x) = \operatorname{mex}\{g(y) : y
\text{ reachable from } x\}$, where $\operatorname{mex}$ is the smallest non-negative integer not in the set. A position with $g(x) = k$
behaves exactly like a Nim heap of size $k$: from it one can move to positions of every Grundy number below $k$, and possibly to some
above $k$, but such a move can be answered by moving back down to $k$. So positions with $g = 0$ are losing, and the
\textbf{Sprague--Grundy theorem} follows: the Grundy number of a sum of games is the XOR of the Grundy numbers of its parts. Every game
built from independent parts is Nim in disguise once the parts' Grundy numbers are known (Figure~\ref{fig:grundy}).

\begin{figure}[htbp]
\centering
\begin{tikzpicture}
  \node[blbl, anchor=east] at (-0.2, 0.75) {heap size $x$};
  \foreach \xv in {0, ..., 14} \node[lbl] at (0.35 + 0.7*\xv, 0.75) {\xv};
  \node[blbl, anchor=east] at (-0.2, 0) {moves $\{1, 2, 3\}$};
  \arrs{Ga}{0}{0}{0/hot, 1/n, 2/n, 3/n, 0/hot, 1/n, 2/n, 3/n, 0/hot, 1/n, 2/n, 3/n, 0/hot, 1/n, 2/n}
  \node[blbl, anchor=east] at (-0.2, -0.9) {moves $\{1, 3, 4\}$};
  \arrs{Gb}{0}{-0.9}{0/hot, 1/n, 0/hot, 1/n, 2/n, 3/n, 2/n, 0/hot, 1/n, 0/hot, 1/n, 2/n, 3/n, 2/n, 0/hot}
\end{tikzpicture}
\caption{Grundy numbers of single heaps for two subtraction games. With moves $\{1, 2, 3\}$, $g(x) = x \bmod 4$; with $\{1, 3, 4\}$ the
values repeat with period 7. Red cells are the losing positions.}
\label{fig:grundy}
\end{figure}

\usetemplate{Grundy numbers}{Ctrl+Alt+Y, G}{cpgrundy}

\begin{problem}{Nim Game II (CSES)}
Nim, except that a move removes only 1, 2 or 3 sticks from a heap. Who wins (up to $2 \cdot 10^5$ heaps in total, sizes up to
$10^9$)?
\solution A heap of size $x$ has Grundy number $x \bmod 4$ (Figure~\ref{fig:grundy}); XOR them. 1000 tests of 200 heaps take 0.08\,s.
\end{problem}

\begin{code}
void solve() {
    int t;
    cin >> t;
    while (t--) {
        int n;
        cin >> n;
        ll x = 0;
        rep(i, n) {
            ll h;
            cin >> h;
            x ^= h % 4;  // the Grundy number of a heap when 1 to 3 are taken
        }
        print(x ? "first" : "second");
    }
}
\end{code}

\begin{trick}[Find the pattern]
The Grundy numbers of a single heap are usually periodic after a short prefix, or follow a simple formula. Compute them up to a few
thousand with the template, print them, and look; then answer heaps of size up to $10^{18}$ from the pattern. Verify the pattern
against the brute force further out than you expect to need.
\end{trick}

\paragraph{Staircase Nim.} Balls lie on stairs $1, \dots, n$, and a move takes any positive number of balls from stair $k \ge 2$ down
to stair $k - 1$; balls on stair 1 can no longer move. Only the even stairs matter: the game is Nim on the even stairs
(Figure~\ref{fig:stairnim}). Moving balls from an even stair to an odd one removes them from that Nim heap; moving balls from an odd stair
to an even one adds to a heap, but the opponent can move exactly those balls one stair further, back onto an odd stair, restoring the
position. So the winner of Nim on the even stairs wins the staircase.

\begin{figure}[htbp]
\centering
\begin{tikzpicture}
  \foreach \stp/\cnt in {1/2, 2/3, 3/1, 4/4, 5/2} {
    \pgfmathtruncatemacro{\ev}{mod(\stp, 2) == 0 ? 1 : 0}
    \ifnum\ev=1
      \fill[fgsky!35] (1.3*\stp - 1.3, 0) rectangle (1.3*\stp, 0.45*\stp);
    \else
      \fill[fgpaper] (1.3*\stp - 1.3, 0) rectangle (1.3*\stp, 0.45*\stp);
    \fi
    \draw[line width=0.8pt, draw=fgink!70] (1.3*\stp - 1.3, 0) rectangle (1.3*\stp, 0.45*\stp);
    \foreach \bl in {1, ..., \cnt}
      \node[circle, draw=fgorange!85!black, fill=fgorange!30, line width=0.6pt, minimum size=3.6mm, inner sep=0pt] at
        (1.3*\stp - 1.3 + 0.27*\bl, 0.45*\stp + 0.2) {};
    \node[lbl] at (1.3*\stp - 0.65, -0.3) {stair \stp};
  }
  \draw[->, line width=0.9pt, draw=fgverm] (4.35, 2.42) to[bend right=40] node[above, font=\sffamily\scriptsize, text=fgverm]
    {a move} (3.45, 1.9);
  \node[notel, anchor=west, text width=5.8cm] at (7.0, 1.2) {Balls 2, 3, 1, 4, 2 on stairs 1 to 5. The even stairs (blue) hold 3 and
    4 balls, and $3 \oplus 4 = 7 \ne 0$, so the first player wins: moving one ball from stair 4 to stair 3 (arrow) leaves $3 \oplus
    3 = 0$ on the even stairs.};
\end{tikzpicture}
\caption{Staircase Nim is Nim on the even stairs.}
\label{fig:stairnim}
\end{figure}

\begin{problem}{Stair Game (CSES)}
Decide the staircase game above for many tests.
\solution XOR the ball counts on the even stairs. 0.08\,s for 1000 tests of 200 stairs.
\end{problem}

\begin{code}
void solve() {
    int t;
    cin >> t;
    while (t--) {
        int n;
        cin >> n;
        ll x = 0;  // staircase Nim: only the balls on even stairs matter
        rep(k, 1, n + 1) {
            ll p;
            cin >> p;
            if (k % 2 == 0) x ^= p;
        }
        print(x ? "first" : "second");
    }
}
\end{code}

\begin{problem}{Another Game (CSES)}
Several heaps of coins; a move takes one coin from each heap in any non-empty set of non-empty heaps. Who takes the last coin?
\solution Not every game needs Grundy numbers; sometimes an invariant suffices. If all heaps are even, any move makes some heap odd,
and the reply ``one coin from every odd heap'' makes them all even again; the final position, all zeros, is all-even. So the player to
move loses exactly when every heap is even.
\end{problem}

\begin{code}
void solve() {
    int t;
    cin >> t;
    while (t--) {
        int n;
        cin >> n;
        bool odd = false;  // all heaps even: the first player loses
        rep(i, n) {
            ll x;
            cin >> x;
            if (x % 2) odd = true;
        }
        print(odd ? "first" : "second");
    }
}
\end{code}

\begin{recognize}
Two players alternate, both have the same moves, there is no chance and no hidden information, and the player who cannot move loses:
classify single games with a DP, combine independent games with Grundy numbers and XOR, and look for invariants (parity, symmetry,
copying the opponent's move) before computing anything. When the player who makes the \emph{last} move loses (mis\`ere play), Nim
changes only when every heap has at most one stone; then the player to move wins exactly when the number of non-empty heaps is even.
\end{recognize}

\begin{summary}[Chapter summary]
\begin{itemize}
  \item Product, sum and complement rules, bijections and double counting build every count; multinomials divide out repeated letters.
  \item Factorials and inverse factorials give $\binom{n}{k}$ in $\Oh{1}$; Pascal's rule works for any modulus; Lucas and Kummer
        handle small primes.
  \item Stars and bars distributes identical items; lattice paths avoiding obstacles are counted by charging each bad path to its first
        obstacle; the reflection principle gives Catalan numbers and ballot counts.
  \item Inclusion--exclusion counts unions, derangements and surjections; Burnside's lemma averages fixed points to count up to symmetry.
  \item Stirling numbers convert powers to falling factorials; Lagrange interpolation evaluates any polynomial from enough values;
        determinants count spanning trees and non-intersecting paths.
  \item Expectations are sums of indicator probabilities; maxima use tail sums.
  \item Nim is decided by XOR; Sprague--Grundy reduces sums of impartial games to Nim; staircase Nim is Nim on the even stairs.
\end{itemize}
\end{summary}

\begin{practice}
\ppgroup{Counting}
\begin{pplist}
  \item Binomial Coefficients
  \item Creating Strings II
  \item Distributing Apples
  \item Christmas Party
  \item Bracket Sequences I
  \item Bracket Sequences II
  \item Counting Necklaces
  \item Counting Grids
  \item Grid Paths II
\end{pplist}
\ppgroup{Probability and games}
\begin{pplist}
  \item Dice Probability
  \item Moving Robots
  \item Candy Lottery
  \item Inversion Probability
  \item Stick Game
  \item Nim Game I
  \item Nim Game II
  \item Stair Game
  \item Grundy's Game
  \item Another Game
\end{pplist}
\end{practice}
